\documentclass[12pt,letterpaper]{article}

% ---------- Paquetes ----------
\usepackage[spanish,es-tabla]{babel}
\usepackage[utf8]{inputenc}
\usepackage[T1]{fontenc}
\usepackage{lmodern}
\usepackage[margin=2.5cm]{geometry}
\usepackage{graphicx}
\usepackage{float}
\usepackage{setspace}
\usepackage{booktabs}
\usepackage{xcolor}
\usepackage{listings}
\usepackage{hyperref}
\hypersetup{colorlinks=true, linkcolor=black, urlcolor=blue, citecolor=black}

\onehalfspacing
\setlength{\parindent}{0pt}
\setlength{\parskip}{0.6em}

\lstset{
  basicstyle=\ttfamily\small,
  backgroundcolor=\color{gray!10},
  frame=single,
  breaklines=true,
  columns=fullflexible
}

% \captura{archivo}{pie de figura}{etiqueta}
% Si la imagen todavía no existe en figuras/, se muestra un recuadro con el nombre esperado.
\newcommand{\captura}[3]{%
  \begin{figure}[H]
    \centering
    \IfFileExists{figuras/#1}{%
      \includegraphics[width=0.92\textwidth,height=0.42\textheight,keepaspectratio]{figuras/#1}%
    }{%
      \fbox{\parbox[c][5cm][c]{0.85\textwidth}{\centering\small Pegar captura:\\ \texttt{figuras/#1}}}%
    }
    \caption{#2}
    \label{#3}
  \end{figure}%
}

% ============================================================
\begin{document}

% ---------- Portada institucional ----------
\begin{titlepage}
  \centering
  \IfFileExists{figuras/logo.png}{\includegraphics[width=4cm]{figuras/logo.png}\par}{}
  \vspace{1cm}
  {\Large\bfseries NOMBRE DE LA UNIVERSIDAD\par}
  \vspace{0.3cm}
  {\large Facultad / División de ...\par}
  {\large Carrera: ...\par}
  \vspace{2cm}
  {\large Gestión del Proceso de Desarrollo\par}
  \vspace{0.8cm}
  {\LARGE\bfseries Proyecto Integrador:\\[0.3cm]
   Pipeline CI/CD automatizado para una API REST\\ con Docker, GitHub Actions y AWS EC2\par}
  \vspace{2cm}
  {\large\bfseries Integrantes:\par}
  {\large Nombre Apellido -- Matrícula\par}
  {\large Nombre Apellido -- Matrícula\par}
  \vspace{1cm}
  {\large\bfseries Profesor(a):\par}
  {\large Nombre del profesor(a)\par}
  \vfill
  {\large Ciudad, \today\par}
\end{titlepage}

\tableofcontents
\listoffigures
\newpage

% ============================================================
\section{Introducción}

En el desarrollo de software actual, escribir código que funcione es solo una parte del trabajo. Igual de importante es la forma en la que ese código llega de la computadora del desarrollador a los usuarios finales: cómo se verifica, cómo se empaqueta y cómo se publica en un servidor. Cuando estas tareas se realizan de manera manual, son lentas, dependen de la memoria de las personas y es fácil cometer errores, como olvidar ejecutar las pruebas, copiar una versión equivocada al servidor o dejar una configuración distinta entre el equipo de desarrollo y el de producción. Para resolver estos problemas surgieron las prácticas de \textit{Integración Continua} (CI) y \textit{Despliegue Continuo} (CD), que forman parte de la cultura DevOps.

La Integración Continua consiste en que cada cambio que se sube al repositorio se compila y se prueba automáticamente, de modo que los errores se detectan en minutos y no semanas después. El Despliegue Continuo extiende esta idea: si las pruebas son exitosas, el mismo sistema empaqueta la aplicación y la publica en el servidor sin intervención humana. De esta forma cada \texttt{git push} se convierte en una nueva versión verificada y disponible para los usuarios, y el equipo puede entregar mejoras pequeñas con mayor frecuencia y menor riesgo \cite{humble2010}.

Para que este flujo sea confiable, la aplicación debe ejecutarse igual en cualquier máquina. Esto se logra con contenedores: Docker permite empaquetar la aplicación junto con su entorno de ejecución (sistema base, versión de Node.js y dependencias) en una \textit{imagen} inmutable \cite{dockerdocs}. Esa imagen se guarda en un registro, en este caso Docker Hub, y desde ahí puede descargarse y ejecutarse en cualquier servidor que tenga Docker instalado. Así se elimina el clásico problema de ``en mi computadora sí funciona''.

En este proyecto se tomó una API REST desarrollada previamente con Node.js, Express y SQLite, que administra un catálogo de libros, autores y categorías, y se construyó a su alrededor un flujo de trabajo profesional y completamente automatizado. El flujo se define en GitHub Actions \cite{githubactions} y está compuesto por tres etapas: primero se ejecuta la batería de pruebas automatizadas con Jest y Supertest \cite{jest}, exigiendo una cobertura mínima de código del 70\%; después se construye la imagen Docker y se publica en Docker Hub con dos etiquetas, \texttt{latest} y el hash del commit; finalmente, el pipeline se conecta por SSH a una instancia de AWS EC2 con Ubuntu Server \cite{awsec2}, descarga la versión más reciente, detiene el contenedor anterior y levanta el nuevo en el puerto 80.

Un aspecto central del proyecto es la seguridad. Ningún dato sensible, como contraseñas, tokens de acceso, direcciones IP o llaves SSH, aparece en el código fuente del repositorio. Todos estos valores se almacenan cifrados en GitHub Secrets y solo se inyectan en el pipeline durante su ejecución. Además, la autenticación con Docker Hub se realiza mediante un \textit{Personal Access Token} en lugar de la contraseña de la cuenta, y el contenedor se ejecuta con un usuario sin privilegios de administrador.

El presente reporte describe la arquitectura de la solución, los pasos de configuración realizados y los resultados obtenidos, incluyendo la evidencia de las pruebas automatizadas, la publicación de la imagen y el despliegue en la nube. Por último, se presenta una demostración en la que un cambio en el código se propaga automáticamente hasta el servidor público con un solo \texttt{git push}.

% ============================================================
\section{Arquitectura de la solución}

La solución está formada por los siguientes componentes:

\begin{itemize}
  \item \textbf{API REST:} Node.js 22, Express 5 y SQLite (\texttt{better-sqlite3}). Expone 11 endpoints para categorías, autores, libros y administración de la base de datos. Todas las respuestas usan el formato \texttt{\{ "statusCode": 200, "data": [] \}}.
  \item \textbf{Pruebas:} Jest y Supertest. Las pruebas usan una base de datos en memoria para no afectar los datos reales.
  \item \textbf{Contenedor:} \texttt{Dockerfile} de dos etapas basado en \texttt{node:22-slim}, ejecutado con un usuario sin privilegios y con un \texttt{HEALTHCHECK}. El archivo \texttt{.dockerignore} excluye \texttt{node\_modules}, \texttt{.env}, logs, pruebas y bases de datos locales.
  \item \textbf{Pipeline:} \texttt{.github/workflows/main.yml} con tres jobs encadenados: \texttt{test}, \texttt{build-and-push} y \texttt{deploy}.
  \item \textbf{Registro:} repositorio \texttt{webapp} en Docker Hub.
  \item \textbf{Servidor:} instancia AWS EC2 con Ubuntu Server y Docker; el Security Group solo permite los puertos 22 (SSH) y 80 (HTTP).
\end{itemize}

El flujo completo es: \textit{git push} $\rightarrow$ pruebas y cobertura $\rightarrow$ construcción y publicación de la imagen $\rightarrow$ conexión SSH a EC2 $\rightarrow$ \texttt{docker pull} $\rightarrow$ reemplazo del contenedor en el puerto 80. Las pruebas se ejecutan también en cada \textit{pull request}, mientras que la publicación y el despliegue solo ocurren con un \textit{push} a la rama \texttt{main}.

\begin{table}[H]
  \centering
  \caption{Endpoints de la API}
  \begin{tabular}{rlll}
    \toprule
    \# & Método & Ruta & Descripción \\
    \midrule
    1  & GET    & /api/categorias        & Listar categorías \\
    2  & POST   & /api/categorias        & Crear categoría \\
    3  & GET    & /api/autores           & Listar autores \\
    4  & POST   & /api/autores           & Crear autor \\
    5  & GET    & /api/libros            & Listar libros \\
    6  & GET    & /api/libros/:id        & Obtener un libro \\
    7  & POST   & /api/libros            & Crear libro \\
    8  & PUT    & /api/libros/:id        & Actualizar libro \\
    9  & DELETE & /api/libros/:id        & Eliminar libro \\
    10 & POST   & /api/database/backup   & Respaldo de la base de datos \\
    11 & DELETE & /api/database/vaciar   & Vaciar la base de datos \\
    \bottomrule
  \end{tabular}
\end{table}

\begin{table}[H]
  \centering
  \caption{GitHub Secrets utilizados por el pipeline}
  \begin{tabular}{ll}
    \toprule
    Secret & Contenido \\
    \midrule
    \texttt{DOCKERHUB\_USERNAME} & Usuario de Docker Hub \\
    \texttt{DOCKERHUB\_TOKEN}    & Personal Access Token de Docker Hub \\
    \texttt{EC2\_HOST}           & IP pública de la instancia \\
    \texttt{EC2\_USER}           & Usuario SSH (\texttt{ubuntu}) \\
    \texttt{EC2\_SSH\_KEY}       & Contenido de la llave \texttt{.pem} \\
    \bottomrule
  \end{tabular}
\end{table}

% ============================================================
\section{Resultados}

\subsection{Pruebas automatizadas y cobertura local}

Antes de automatizar el flujo se verificó que la batería de pruebas funcionara en el equipo de desarrollo con el comando \texttt{npm run test:coverage}. Se ejecutaron 51 pruebas de integración distribuidas en 5 archivos, que cubren los casos exitosos y los casos de error de cada endpoint (datos faltantes, formatos inválidos, registros duplicados, identificadores inexistentes, JSON mal formado y fallos internos). Todas las pruebas pasaron y la cobertura obtenida fue de 96.59\% en sentencias, 89.89\% en ramas, 100\% en funciones y 97.45\% en líneas, por encima del umbral mínimo de 70\% configurado en Jest (Figura~\ref{fig:tests}).

\captura{fig01-pruebas-local.png}{Ejecución local de las pruebas con reporte de cobertura}{fig:tests}

\subsection{Construcción y ejecución de la imagen Docker}

Con el comando \texttt{docker build -t practica3-api .} se construyó la imagen. La primera etapa instala únicamente las dependencias de producción y la segunda copia solo el código necesario, por lo que la imagen final no contiene pruebas, archivos de configuración local ni dependencias de desarrollo (Figura~\ref{fig:build}).

\captura{fig02-docker-build.png}{Construcción de la imagen Docker}{fig:build}

Después se ejecutó el contenedor localmente con \texttt{docker run}, mapeando el puerto 8080 del equipo al 3000 del contenedor. El comando \texttt{docker ps} muestra el estado \texttt{healthy}, lo que indica que el \texttt{HEALTHCHECK} verificó que la API responde (Figura~\ref{fig:run}). Al consultar los endpoints desde el cliente HTTP se obtuvieron las respuestas esperadas (Figura~\ref{fig:apilocal}).

\captura{fig03-docker-run.png}{Contenedor en ejecución con estado \textit{healthy}}{fig:run}
\captura{fig04-api-local.png}{Respuesta de la API ejecutándose dentro del contenedor}{fig:apilocal}

\subsection{Configuración de Docker Hub}

En Docker Hub se utilizó el repositorio \texttt{webapp} y se creó un \textit{Personal Access Token} con permisos de lectura y escritura. Este token es el que utiliza el pipeline para autenticarse, de modo que la contraseña de la cuenta nunca se usa ni se almacena (Figura~\ref{fig:token}).

\captura{fig05-dockerhub-token.png}{Personal Access Token creado en Docker Hub}{fig:token}

\subsection{Infraestructura en AWS EC2}

Se lanzó una instancia EC2 con Ubuntu Server (Figura~\ref{fig:ec2}). Su Security Group solo permite tráfico entrante por el puerto 22, para la conexión SSH del pipeline, y por el puerto 80, para las peticiones HTTP a la API (Figura~\ref{fig:sg}). Posteriormente se instaló Docker en el servidor y se agregó el usuario \texttt{ubuntu} al grupo \texttt{docker} (Figura~\ref{fig:ec2docker}).

\captura{fig06-ec2-instancia.png}{Instancia EC2 en ejecución}{fig:ec2}
\captura{fig07-security-group.png}{Reglas de entrada del Security Group (puertos 22 y 80)}{fig:sg}
\captura{fig08-ec2-docker.png}{Docker instalado en la instancia EC2}{fig:ec2docker}

\subsection{GitHub Secrets}

Los cinco valores sensibles se registraron como secretos del repositorio. GitHub los almacena cifrados, no permite volver a leerlos desde la interfaz y los oculta automáticamente en los logs del pipeline (Figura~\ref{fig:secrets}).

\captura{fig09-github-secrets.png}{Secrets configurados en el repositorio}{fig:secrets}

\subsection{Ejecución del pipeline}

Al hacer \texttt{git push} a la rama \texttt{main}, GitHub Actions ejecutó los tres jobs en orden. Cada uno depende del anterior, por lo que si las pruebas fallan no se publica ninguna imagen y el servidor conserva la versión estable (Figura~\ref{fig:pipeline}).

\captura{fig10-actions-pipeline.png}{Pipeline completo en GitHub Actions}{fig:pipeline}

En el job \texttt{test} se muestra la tabla de cobertura en los logs y un resumen en la página de la ejecución (Figura~\ref{fig:cov}).

\captura{fig11-actions-cobertura.png}{Resumen de cobertura en los logs del pipeline}{fig:cov}

El job \texttt{build-and-push} publicó la imagen en Docker Hub con las etiquetas \texttt{latest} y el hash del commit, lo que permite identificar exactamente qué versión del código contiene cada imagen y regresar a una anterior si fuera necesario (Figura~\ref{fig:tags}).

\captura{fig12-dockerhub-tags.png}{Etiquetas \texttt{latest} y hash del commit en Docker Hub}{fig:tags}

El job \texttt{deploy} se conectó por SSH a la EC2, descargó la imagen, reemplazó el contenedor y comprobó que la API respondiera en el puerto 80 (Figura~\ref{fig:deploy}).

\captura{fig13-deploy-log.png}{Log del despliegue en EC2}{fig:deploy}

\subsection{API pública}

Una vez terminado el despliegue, la API quedó disponible en la dirección pública de la instancia (Figura~\ref{fig:publica}).

\captura{fig14-api-ec2.png}{API respondiendo desde la IP pública de EC2}{fig:publica}

\subsection{Demostración de despliegue automático}

Para comprobar el flujo completo se cambió el mensaje de error 404 en \texttt{src/app.js} y se hizo \texttt{git push}. Antes del cambio, una ruta inexistente devolvía el mensaje original (Figura~\ref{fig:antes}); tras unos minutos, sin entrar al servidor, la misma petición devolvió el mensaje nuevo (Figura~\ref{fig:despues}). Esto muestra que el pipeline ejecuta las pruebas, actualiza Docker Hub y refresca el servidor automáticamente.

\captura{fig15-demo-antes.png}{Respuesta antes del cambio}{fig:antes}
\captura{fig16-demo-despues.png}{Respuesta después del \texttt{git push}}{fig:despues}

% ============================================================
\section{Conclusiones}

La implementación del pipeline demostró que automatizar el ciclo de vida de una aplicación reduce de forma importante el trabajo manual y los errores asociados a él. Con un solo \texttt{git push} el código se prueba, se empaqueta en una imagen y se publica en un servidor en la nube en pocos minutos, sin que ninguna persona tenga que conectarse al servidor. Las pruebas automatizadas, junto con el umbral mínimo de cobertura, funcionan como una barrera de calidad: una versión que rompe alguna funcionalidad nunca llega a producción, y el servidor conserva la última versión estable. Además, etiquetar cada imagen con el hash del commit permite saber con exactitud qué código está en ejecución y regresar a una versión anterior si fuera necesario.

Por otra parte, el proyecto mostró la importancia de la seguridad y de la portabilidad en este tipo de flujos. Al manejar contraseñas, tokens, direcciones IP y llaves SSH únicamente mediante GitHub Secrets, el repositorio puede ser público sin exponer la infraestructura. Docker, por su parte, garantizó que la aplicación se comportara igual en la computadora de desarrollo, en los servidores de GitHub y en la instancia EC2. Como mejora futura se podría colocar un proxy inverso o un balanceador de carga delante de los contenedores para lograr despliegues completamente sin interrupción, y agregar etapas de análisis de seguridad de la imagen antes de publicarla.

% ============================================================
\begin{thebibliography}{9}

\bibitem{humble2010}
J. Humble y D. Farley, \textit{Continuous Delivery: Reliable Software Releases through Build, Test, and Deployment Automation}. Addison-Wesley, 2010.

\bibitem{dockerdocs}
Docker Inc., ``Docker Documentation''. [En línea]. Disponible en: \url{https://docs.docker.com/}

\bibitem{githubactions}
GitHub, ``GitHub Actions documentation''. [En línea]. Disponible en: \url{https://docs.github.com/en/actions}

\bibitem{awsec2}
Amazon Web Services, ``Amazon EC2 User Guide''. [En línea]. Disponible en: \url{https://docs.aws.amazon.com/ec2/}

\bibitem{jest}
Jest, ``Jest Documentation -- Configuring Jest (coverageThreshold)''. [En línea]. Disponible en: \url{https://jestjs.io/docs/configuration}

\end{thebibliography}

\end{document}
